\section{Дневник выполнения работы}

\subsection{Подготовка}
Поставил {\ttfamily valgrind}, {\ttfamily libdmalloc-dev} и {\ttfamily time} из репозитория Debian.
Исходник первой лабораторной скопировал в {\ttfamily old/B.cpp}, исправленную версию вёл в
{\ttfamily new/B.cpp}. Написал генератор тестов {\ttfamily gen.py} на Python (режимы {\ttfamily mix},
{\ttfamily insert}, {\ttfamily sorted}, {\ttfamily saveload}) и сгенерировал 5 тестов.

Первый вариант режима {\ttfamily sorted} делал слова вида {\ttfamily aa\ldots aab}, у которых длина росла
до 1500 символов "--- это больше допустимых 256. Переделал на запись номера в 26-ричной системе из 5 букв.

\subsection{Замер времени старой версии}
\begin{alltt}
\$ g++ -O2 -o b old/B.cpp
\$ /usr/bin/time -f '\%es \%MKB' ./b < tests/mix_1000000.txt > /dev/null
mix_1000000 0.89s 31948KB
insert_1000000 1.37s 82536KB
sorted_1000000 0.65s 81456KB
saveload_200000 0.89s 35948KB
\end{alltt}
На $10^6$ вставок уходит 82 МБ, то есть около 82 байт на слово, хотя средняя длина слова около 10 букв.
Это первая подсказка, что узел хранится расточительно.

\subsection{gprof, старая версия}
\begin{alltt}
\$ g++ -O2 -pg -o bpg old/B.cpp
\$ ./bpg < tests/mix_1000000.txt > /dev/null
\$ gprof -b -p bpg gmon.out
  \%   cumulative   self              self     total
 time   seconds   seconds    calls  ns/call  ns/call  name
 25.73      0.09     0.09 11818239     7.62     7.62  balance(Node*)
 25.73      0.18     0.09   300389   299.74   299.74  find(Node*, string const&)
 20.01      0.25     0.07   499745   140.13   262.48  insert(Node*, string const&, ...)
 11.43      0.29     0.04  1000000    40.02    40.02  lower(string)
 11.43      0.33     0.04   199866   200.22   318.13  erase(Node*, string const&, bool&)
\end{alltt}
Сумма по профилю "--- 0.35~с, а программа работает 0.89~с. Значит, больше половины времени уходит
в функции стандартной библиотеки, которых {\ttfamily gprof} не видит. Ещё заметно, что
{\ttfamily balance} вызывается 11.8 млн раз "--- в том числе на обратном пути рекурсии, когда
слово уже было в словаре и дерево не менялось.

Профиль на тесте {\ttfamily saveload\_200000}:
\begin{alltt}
 35.50      0.11     0.11  1981950     0.00     0.00  insert(...)
 25.82      0.19     0.08 32841219     0.00     0.00  balance(Node*)
 12.91      0.23     0.04       10     4.00    22.89  load(Node*&, string const&)
\end{alltt}
Один {\ttfamily Load} стоит 22.9~мс, и почти всё это время "--- {\ttfamily insert} каждого слова
с поворотами, хотя в файле слова уже лежат в известном порядке.

\subsection{callgrind, старая версия}
Чтобы увидеть, куда делась остальная половина времени, запустил {\ttfamily callgrind}:
\begin{alltt}
\$ valgrind --tool=callgrind ./bg < tests/mix_100000.txt > /dev/null
\$ callgrind_annotate --inclusive=yes callgrind.out
489,396,563 (100.0\%)  PROGRAM TOTALS
151,471,892 (30.95\%)  std::basic_istringstream<...>::basic_istringstream(string)
129,243,054 (26.41\%)  std::basic_ios<...>::init(...)
115,113,857 (23.52\%)  std::basic_ios<...>::_M_cache_locale(...)
 94,969,870 (19.41\%)  __dynamic_cast
 59,259,094 (12.11\%)  std::istream::_M_extract<unsigned long>(...)
\end{alltt}
Почти треть всех инструкций "--- это {\ttfamily istringstream}, который создаётся заново на каждую
строку ввода. Конструктор потока каждый раз инициализирует локаль и делает {\ttfamily dynamic\_cast}.
Ещё 12\% "--- разбор числа через поток. Сама работа с деревом оказалась не самой дорогой частью.

\subsection{memcheck и ручные граничные тесты}
\begin{alltt}
\$ valgrind --leak-check=full ./bg < tests/mix_100000.txt > /dev/null
==   total heap usage: 213,834 allocs, 213,828 frees, 8,113,110 bytes allocated
==    definitely lost: 0 bytes in 0 blocks
==    still reachable: 122,880 bytes in 6 blocks
== ERROR SUMMARY: 0 errors from 0 contexts
\end{alltt}
Утечек нет, 6 блоков "--- внутренние буферы {\ttfamily libstdc++}. Но аллокаций 213 тыс., хотя
вставок всего 50 тыс.: на каждую строку {\ttfamily substr} и {\ttfamily istringstream} выделяют память.

Потом написал файл с некорректными командами и прогнал через старую версию:
\begin{alltt}
+ abc                         -> OK         (номера нет, записан мусор)
+ x1y 5                       -> OK         (цифра в слове)
+ bbbb...b (300 букв) 7       -> OK         (длиннее 256)
! Save s.bin                  -> OK
! Load s.bin                  -> ERROR: Wrong file format
! Foo zzz                     -> OK         (словарь очищен!)
! Load                        -> OK         (словарь очищен!)
+ big 99999999999999999999    -> OK, потом big -> OK: 18446744073709551615
\end{alltt}
Получилось, что программа сама сохраняет файл, который потом не может загрузить. Подробно все
недочёты разобраны в следующем разделе.

\subsection{dmalloc, старая версия}
\begin{alltt}
\$ g++ -O2 -g -o bdm old/B.cpp -ldmalloccxx
\$ eval `dmalloc -b -l dm.log -i 100 high`
\$ ./bdm < tests/mix_100000.txt > /dev/null
alloc calls: malloc 7, calloc 0, realloc 0, free 0
alloc calls: new 213827, delete 213827
 total memory allocated: 8113110 bytes (213834 pnts)
 max in use at one time: 2542105 bytes (41245 pnts)
\end{alltt}
Работает в 60 раз медленнее (8.5~с против 0.14~с) из-за проверки кучи. Ошибок заборов нет,
{\ttfamily new} и {\ttfamily delete} парные. Семь {\ttfamily malloc} без {\ttfamily free} "--- те же
буферы {\ttfamily libstdc++} и {\ttfamily stdio}, что видел {\ttfamily memcheck}. Для {\ttfamily new}
dmalloc не пишет файл и строку ({\ttfamily from 'unknown'}), потому что макросы из
{\ttfamily dmalloc.h} работают только для {\ttfamily malloc}.

\subsection{massif, старая версия}
\begin{alltt}
\$ valgrind --tool=massif ./bg < tests/saveload_200000.txt > /dev/null
\$ ms_print massif.out
    MB
29.10^                    #
     |                   @#     @          :     :     :    :     :     :    :
     |                   @#    :@   ::    @:    ::    @:   :@    :@    @:   :@
     |                 @@@#  :::@  :: :  :@:  ::::  ::@:  ::@   ::@  ::@:  ::@
     |              @:@@ @# ::::@ ::: : ::@:  : :: :::@:::::@ ::::@ :::@:::::@
     |         :::@@@:@@ @#:::::@:::: ::::@:: : ::@:::@:::::@:::::@::::@:::::@
     | @:@@@@@@:::@@@:@@ @#:::::@:::: ::::@:: : ::@:::@:::::@:::::@::::@:::::@
   0 +----------------------------------------------------------------------->Gi
     0                                                                   5.870
  n        time(i)         total(B)   useful-heap(B) extra-heap(B)
 24  1,800,565,251       30,515,256       25,987,062     4,528,194
\end{alltt}
Зубцы "--- это {\ttfamily Load}: новое дерево строится рядом со старым, и только при успехе старое
удаляется. Так и должно быть по условию (при ошибке словарь не должен меняться), поэтому пик в два
размера словаря неизбежен, но сам размер словаря можно уменьшить.

\subsection{Исправления}
По результатам переписал {\ttfamily new/B.cpp}:
\begin{enumerate}
\item Узел хранит ключ прямо в себе: {\ttfamily char key[1]} в конце структуры, память выделяется
одним {\ttfamily malloc(offsetof(Node, key) + len)}. Длина хранится как {\ttfamily len - 1} в одном
байте (1--256 помещается в 0--255), высота "--- в {\ttfamily signed char}.
\item Строка разбирается вручную функциями {\ttfamily readWord}, {\ttfamily readNum},
{\ttfamily atEnd} по индексу в той же строке, без {\ttfamily substr} и {\ttfamily istringstream}.
Сразу проверяется формат: только буквы, длина не больше 256, число есть и не переполняет
{\ttfamily uint64\_t}, после аргументов ничего лишнего. На неверную команду выводится
{\ttfamily ERROR: Invalid command}.
\item {\ttfamily insert} и {\ttfamily erase} не вызывают {\ttfamily balance} на обратном пути,
если дерево не изменилось.
\item {\ttfamily Save} пишет узлы в симметричном порядке (по возрастанию), формат {\ttfamily AVL2}.
{\ttfamily Load} читает узлы в массив, проверяя, что ключи строго возрастают, и строит идеально
сбалансированное дерево за $O(n)$ функцией {\ttfamily link} "--- без единого поворота.
\item {\ttfamily Save} пишет во временный файл {\ttfamily path.tmp} и делает {\ttfamily rename}
только после успешной записи.
\end{enumerate}

\subsection{Проверка исправленной версии}
Сначала пример из условия и граничные тесты под {\ttfamily -fsanitize=address,undefined}. Потом
стресс-тест: 5 прогонов по 30000 случайных команд, включая {\ttfamily Save}/{\ttfamily Load},
сравнение с эталоном на {\ttfamily dict} в Python. Все совпали.

Подключил {\ttfamily dmalloc.h} через {\ttfamily \#ifdef DMALLOC}, как в методичке. Так как узлы
теперь выделяются {\ttfamily malloc}, dmalloc стал показывать строку:
\begin{alltt}
\$ g++ -O2 -g -DDMALLOC -DDMALLOC_FUNC_CHECK -o bdm new/B.cpp -ldmalloccxx
top 10 allocations:
 total-size  count in-use-size  count  source
         54      2           0      0  B.cpp:33
\end{alltt}
Для проверки, что dmalloc действительно ловит утечки, собрал копию без {\ttfamily clear(root)} в
конце {\ttfamily main}:
\begin{alltt}
Dumping Not-Freed Pointers Changed Since Start:
 not freed: '0x7f94b4902f88|s1' (27 bytes) from 'leak.cpp:33'
\end{alltt}
27 байт = 26 байт заголовка узла + 1 буква: не освободился узел со словом {\ttfamily b}.

Покрытие {\ttfamily gcov} после всех тестов "--- 92\% строк. Не покрыты только ветки ошибок записи и
нехватки памяти, их проверил отдельно через {\ttfamily ulimit} (см. раздел сравнения).
\pagebreak
